\section{The format}
\label{sec:format}

This section describes the wire format at the level the evaluation needs. The
byte-level specification, and the measurements behind each decision, are in
the library's repository (Section~\ref{sec:availability}).

\subsection{Messages are key runs}

A message is a \emph{key run}: a sequence of fields, each a key, a descriptor
and a payload, ended by the length of whatever encloses it. A run uses one of
two key widths, chosen per struct type and recorded where the run opens. In
the narrow width a field header is a single byte, four bits of key and four
of descriptor, which covers types whose field numbers are 1 to 16. In the wide
width it is two bytes, eight bits each, for numbers up to 255; a type with
more fields is split into linked pages of 255. A nested struct opens its own
run, so narrow and wide types nest in either order. In both widths the
descriptor is enough to step over a field the reader does not know, so
producer and consumer can add and drop fields independently.

Field numbers come from a struct tag (\code{cb:"3"}) or, without one, from an
8-bit hash of the field's name with linear probing on collision. Untagged
types therefore always use the wide width and cost one more byte per field
present, and are limited to 255 fields. Tags are optional; every type in the
evaluation has them.

\subsection{Scalars}

A field holding its type's zero value is not written. An integer is written as
sign and magnitude in the fewest bytes that hold it, one to eight, with the
byte count in the descriptor; in the narrow width, the values the descriptor
itself can name (1 to 4, or 1 to 3 and $-1$ for a signed field) need no
payload, so a boolean or a small enumeration is one byte with its key. A
string of up to eight bytes takes its size from the descriptor; a longer one
carries a length of one, three or five bytes.

No size anywhere in the format is a variable-length integer. Protocol
Buffers' varints are compact, but each is decoded by a loop whose trip count
is data, a cost that byte-oriented schemes such as Stream
VByte~\cite{lemire2018streamvbyte} exist to avoid. In \colbin{} every read of
a size is one load of a width the header has already given.

A scalar floating-point field is written as its IEEE~754 bit pattern with the
bytes reversed, in the unsigned integer form. The zero bytes of a float are
its low mantissa bytes, where an integer's are its high ones, so the trim that
shortens an integer shortens a round float too: \code{1.0} is two bytes.

\subsection{Composites, lists and tables}

A nested struct, a list, a table or a map is written with its byte length in
front, patched in after the body is written, which is what lets a reader skip
it whole. A nil pointer is omitted; a pointer to a zero value is the one case
in which the format writes a zero, to keep the two apart.

A slice of structs is written in one of two layouts
(Figure~\ref{fig:layouts}):
\begin{itemize}
\item as a \emph{list}, each element its own key run, when it has fewer than
  eight elements, or when the element type holds something that is not a
  column (a map, a nested struct, a slice of scalars);
\item as a \emph{table}, when it has at least eight elements and every field
  of the element type is a boolean, an integer, a float, a string, a pointer
  to one of those, or itself a slice of structs that qualifies.
\end{itemize}
The decision is taken per value at encode time, and the two layouts have
different descriptors, so a reader dispatches on what it finds and neither
side is configured. A table writes one key per column instead of one per field
per row. A nullable column (a pointer field, which is how a Go type says a value
may be absent) carries one bit of presence per row and the values of the rows
present only, so the column codec sees real values and a run of nulls costs a
bit a row; a column whose values are all zero, or all null, is not written at
all. A table of more than 4\,096 rows is
written in groups of 4\,096 rows, each group a complete set of columns, which
bounds the memory the encoder transposes through and keeps a group's columns
in cache. The threshold of eight rows is the subject of the first ablation in
Section~\ref{sec:ablations}.

\begin{figure}[t]
\centering\small
\begin{tabular}{@{}l@{\qquad}l@{}}
\textbf{list} (fewer than 8 rows) & \textbf{table} (8 rows or more)\\[2pt]
\code{[len]} & \code{[len][rows]}\\
\code{\ \ [len] k1 v k2 v k3 v} & \code{\ \ k1 [column: v v v v ...]}\\
\code{\ \ [len] k1 v k2 v k3 v} & \code{\ \ k2 [column: v v v v ...]}\\
\code{\ \ [len] k1 v k2 v k3 v} & \code{\ \ k3 [column: v v v v ...]}\\
\end{tabular}
\caption{The two layouts of a slice of structs with three fields. A list
repeats each key in every row; a table writes each key once and its values
together, where a column transform can use what they have in common.}
\label{fig:layouts}
\end{figure}

\subsection{The column codec}
\label{sec:column-codec}

An integer column (and a boolean, date or timestamp column, which are integer
columns) is one transform followed by blocks of bit-packed residuals:
\begin{itemize}
\item \emph{raw}: the values, zigzag-encoded if any is negative;
\item \emph{delta}: differences between consecutive values, zigzag-encoded,
  after the first value;
\item \emph{frame of reference}: each value minus the column's
  minimum~\cite{zukowski2006pfor};
\item \emph{constant}: the single value, and no blocks.
\end{itemize}
The encoder scores each candidate by the bytes its blocks would occupy and
keeps the smallest. Residuals are then packed in blocks of 128 values, each
block at its own bit width, written in one byte ahead of it. The block size is
that of SIMD-BP128~\cite{lemire2015simdbp128}, and it is chosen for the same
arithmetic: 128 values of $w$ bits are exactly $16w$ bytes, so every block
starts and ends on a byte boundary, no state crosses from one block to the
next, and one outlier widens 128 values rather than the column. A block of
zeros has width 0 and costs one byte. Unpacking is one unaligned 8-byte load,
a shift and a mask per value for widths up to 57 bits, with the width switched
on once per block; there is no patching of exceptions as in PFOR. The second
ablation in Section~\ref{sec:ablations} forces each transform in turn.

A column of 64-bit floats whose values are decimals, as most floats that
came out of a database or a CSV file are, is written as integers: the encoder
looks for one exponent $e$ per column such that $v \cdot 10^e$ is an integer
that converts back to $v$ exactly, writes those integers through the integer
codec, and keeps the values that do not fit as exceptions with their position
and bit pattern. This is the decimal half of ALP~\cite{afroozeh2023alp} with
one exponent per column and a division on decode rather than a multiplication
by the inverse, which is slower but exact on every platform. The column is
written this way only when it is smaller than the bit patterns would be.

String columns are written as the strings one after another, each with its
size; they are not dictionary-coded.

\subsection{Optional parts}

\paragraph{The schema section.} The wire carries no names and no types. A
\emph{schema section}, the type's encoding plan serialized as a key, a name
and an operation code per field, can be sent once per connection, or embedded
at the front of a message. With it embedded, the message is self-describing:
it can be rendered as JSON, or decoded into a different type by matching
field names, with no knowledge of the producer's type. The evaluation
measures this mode as \fmt{colbin-json}, and reports the size of the section
against the message in Section~\ref{sec:composition}.

\paragraph{Dynamic values.} A field of an open type (Go's \code{any},
\code{[]any}, \code{map[string]any}) is written with a type descriptor per
value, as schemaless formats do. No case of the evaluation uses this path.

\paragraph{packed5.} An opt-in encoding packs strings of letters and digits
at about five bits per character, when that is shorter. It is off by default
and reported as a separate variant, \fmt{colbin+packed5}.

\subsection{Schema evolution}
\label{sec:evolution}

The writer's type and the reader's need not be the same type. A field is
found by its key, and a key the reader does not know is stepped over by its
descriptor, so a field can be added or removed on either side: the reader
of an older message finds the new field at its zero value, and the reader of
a newer one skips what it does not know. The same holds for a column of a
table, which is one key for all its rows. With tags the key is the tag, so a
field can be renamed freely; without them the key is a hash of the name, and
a renamed field is another field. An integer is written by its value, not by
its declared width, so a field can be widened (an \code{int32} read as an
\code{int64}) or narrowed, as long as the values fit.

Table~\ref{tab:evolution} sets this beside the two formats the evaluation
compares \colbin{} with most closely. It is close to Protocol Buffers'
rules~\cite{protobuf}, with tags in the place of field numbers, and with two
differences against \colbin{}. A value that does not fit the reader's
narrower integer is an error in a record but is truncated without one in a
column, where the check is not made per value. And a change of kind, an
integer to a float or to a string, is not supported and is not always
detected: the reader may return a wrong value instead of an error.
Protocol Buffers also keeps the unknown fields of a message it decodes and
writes again; \colbin{} drops them. The behaviour of \colbin{} in the table
is checked by a test in the artifact (\code{formats/evolution\_test.go});
that of the other two is their documentation's, and for MessagePack's
integers the same test.

\begin{table}[t]
\centering\footnotesize
\begin{tabular}{@{}p{0.23\linewidth}p{0.25\linewidth}p{0.21\linewidth}p{0.21\linewidth}@{}}
\toprule
change & \colbin{} & Protocol Buffers & MessagePack (generated, keyed by name) \\
\midrule
a field or a column added or removed & compatible both ways & compatible both ways & compatible both ways \\
a field renamed & compatible with tags; another field without & compatible & another field, unless a tag keeps the name \\
an integer widened & compatible & compatible & compatible \\
an integer narrowed, a value out of range & error in a record; truncated in a column & truncated & error \\
a change of kind (integer to float or string) & not supported, not always detected & not compatible & error \\
\bottomrule
\end{tabular}
\caption{What a reader sees when its type is not the writer's. \emph{Compatible
both ways}: an old reader decodes a new message and a new reader an old one,
with the missing field at its zero value.}
\label{tab:evolution}
\end{table}

\subsection{Implementation}

The Go implementation has no dependencies. It builds an encoding plan per
type on first use, by reflection, and then runs it over unsafe pointers, so
the steady state involves no reflection except for maps. The hot loops of the
column codec have assembly kernels on amd64 and arm64. Decoders treat their
input as hostile: every length, count and row total is bounded before
anything is allocated. Implementations in Rust and in JavaScript (through
WebAssembly) read and write the same bytes; this paper evaluates the Go
implementation only.
